% context: teacher solutions for 4.1 Mechanical Advantage problem set
\documentclass[12pt, twoside]{article}
\input{../preamble}
\usepackage{float}

\title{Physics}
\author{Chris Huson}
\date{May 2026}

\fancyhead[LE]{\thepage}
\fancyhead[RO]{\thepage}
\fancyhead[LO]{La Scuola d'Italia / Huson / Physics --- Solutions \\* 29 May 2026}

\begin{document}

\subsubsection*{3.6 Problem Set: Mechanical Advantage --- Solutions}

\begin{enumerate}[itemsep=0.4cm]

\item \textbf{Pulleys.}
  \begin{enumerate}[itemsep=0.25cm]
    \item MA $=1$. A fixed pulley does not multiply force; it only changes the \emph{direction} of the effort (pull down to lift up).
    \item Two supporting strands $\Rightarrow$ MA $=2$, so $F_\text{in}=\dfrac{240}{2}=\mathbf{120~\text{N}}$.
    \item MA $=4$: $F_\text{in}=\dfrac{600}{4}=\mathbf{150~\text{N}}$. Rope pulled $=4\times0.50=\mathbf{2.0~\text{m}}$ (force divided by 4, distance multiplied by 4).
  \end{enumerate}

\item \textbf{Lever.}
  \begin{enumerate}[itemsep=0.25cm]
    \item MA $=\dfrac{\text{effort arm}}{\text{load arm}}=\dfrac{1.20}{0.20}=\mathbf{6}$.
    \item $F_\text{in}=\dfrac{F_\text{out}}{\text{MA}}=\dfrac{850}{6}\approx\mathbf{142~\text{N}}$.
    \item Rock rises $\dfrac{0.30}{6}=0.05~\text{m}=\mathbf{5~\text{cm}}$. Trade-off: you gain force ($6\times$) but lose distance ($6\times$) --- work in $=$ work out.
  \end{enumerate}

\item \textbf{Chain and stuck car.}
  \begin{enumerate}[itemsep=0.25cm]
    \item $T=\dfrac{F_\text{pull}}{2\sin\theta}=\dfrac{300}{2\sin10^\circ}=\dfrac{300}{2(0.1736)}\approx\mathbf{864~\text{N}}$.
    \item $\text{MA}=\dfrac{1}{2\sin10^\circ}=\dfrac{1}{0.347}\approx\mathbf{2.88}$ (matches $864/300$).
    \item As $\theta\to0$, $\sin\theta\to0$ and the force on the car $\to\infty$ --- a small sideways pull on a nearly straight chain produces an enormous tension. Practical danger: the chain (or attachment) can snap and whip back; the force can far exceed the chain's rating.
  \end{enumerate}

\item \textbf{Hydraulic lift.}
  \begin{enumerate}[itemsep=0.25cm]
    \item Small: $r=1.0~\text{cm}$, $A=\pi(1.0)^2\approx\mathbf{3.14~\text{cm}^2}$. \;
          Large: $r=15~\text{cm}$, $A=\pi(15)^2\approx\mathbf{707~\text{cm}^2}$.
    \item $\text{MA}=\dfrac{A_\text{large}}{A_\text{small}}=\dfrac{707}{3.14}=\mathbf{225}$ (equivalently $(15/1)^2=225$).
    \item $F_\text{out}=225\times50=\mathbf{1.13\times10^4~\text{N}}$ ($11{,}250~\text{N}$, $\approx$ the weight of an $1150~\text{kg}$ car).
    \item $1500~\text{PSI}\times\dfrac{6.895~\text{kPa}}{1~\text{PSI}}=\mathbf{1.03\times10^4~\text{kPa}}$ ($10{,}340~\text{kPa}=10.3~\text{MPa}$).
  \end{enumerate}

\end{enumerate}

\end{document}
